\section{Introduction}\label{sec:intro}
Teams of unmanned aerial vehicles (UAVs) are used for inspection, delivery and monitoring. During a mission, new tasks
appear, vehicles fail and deadlines are moved forward, and a central planner has to repair the routes of the whole fleet
within a fraction of a second. Multi-robot task allocation and dynamic vehicle routing offer many methods for this
repair~\cite{gerkey2004formal,psaraftis2016dynamic,soeffker2022review,skaltsis2023review}. Auction methods such as the
consensus-based bundle algorithm (CBBA)~\cite{choi2009consensus} and metaheuristics such as adaptive large
neighbourhood search (ALNS)~\cite{ropke2006adaptive} find good plans within such short budgets.

Most of these replanners plan for nominal execution. When a UAV meets wind, flies slower than commanded or stays longer
at a task, the slack of its plan is used up by errors that the planner does not see, and the tasks near the end of a
route are served late. Stochastic and robust routing account for such
errors~\cite{agra2013robust,ehmke2015ensuring,rostami2021correlated,cheng2024weather}, and chance-constrained
allocation bounds the probability of failure~\cite{ponda2012chance,gao2025drcc}. These models are usually solved by
sampling or exact optimisation, which takes longer than the time between two events. A time buffer in the
schedule~\cite{jabali2015selfimposed} is fast, but a constant buffer treats every task alike, although the error grows
with the time until a task is reached. Work on anytime
planning treats the computation time itself as the quantity to
control~\cite{zilberstein1996using,svegliato2018meta,herrmann2023metareasoning}, and it is not clear whether the
planning model or the control of computation matters more when execution is uncertain. A further question concerns the
tasks that a chance constraint gives up. Whether such a task is still worth attempting depends on what a late arrival
costs. Soft time windows price lateness inside the search~\cite{tas2013stochastic}, and a buffer cannot answer the
question at all, because it carries no probability of meeting the deadline.

This paper proposes \mname{} (lead-time-scaled chance constraints), a replanner that keeps the fast heuristics of the
nominal model and changes only the deadlines that they check. On development flights, the spread of the finish-time
error was well approximated by an affine function of the lead time, consistent with wind and airspeed errors that act
over a whole leg. With this error model, a chance constraint on the finish time of every task becomes an effective deadline in closed form. A cost-aware fallback
uses the probability of meeting the deadline to decide which of the given-up tasks are still worth attempting, and it
also takes the form of one deadline per task.

The contribution is threefold. The first is the planning model. A chance constraint on every task reduces exactly to an
effective deadline (Proposition~\ref{prop:cc}), and the cost-aware fallback attempts a given-up task only when its
expected contribution is not negative (Proposition~\ref{prop:costfb}). It reuses the safety factor and the error model
of the planning model and adds no parameter that has to be tuned for the cost of a late arrival, whereas a time buffer
has to be tuned again for every such cost. Gaussian chance
constraints, feasibility checks under uncertainty inside routing
heuristics~\cite{ehmke2015ensuring,bomboi2022stochastic}, ALNS and CBBA are not new, and the contribution lies in the
interface between them and in the decision rule built on it. The second is an evaluation in seven studies on held-out missions, each with a protocol fixed before
its runs. It compares \mname{} with re-implemented baselines from the literature, namely time buffers, moment
propagation~\cite{ehmke2015ensuring}, stopping rules from
metareasoning~\cite{boddy1994deliberation,zilberstein1996using,svegliato2018meta} and CBBA~\cite{choi2009consensus}, and
it separates the effects of the planning model, the fallback, the stopping rule and the trigger. The third is a
software-in-the-loop evaluation with PX4 autopilots~\cite{meier2015px4} in Gazebo~\cite{koenig2004gazebo} and
ROS~2~\cite{macenski2022ros2}, on missions rescaled so that the decision problem stays the same.

In the kinematic tier, \mname{} served more reward than the nominal model, a constant buffer and a conventional planner
on three independent sets of held-out seeds, and \val{WURelCACBBA}\% more than CBBA. A tuned proportional buffer and
moment propagation, whose margins also grow with the lead time, reached a similar utility, so the shape of the margin
carries most of the gain. When a late arrival cost one task reward, the cost-aware fallback scored above the nominal
fallback and above competitors tuned for that cost. In software-in-the-loop flights, the fallback raised the utility by
\val{HURelNoFb}\%, and the kinematic tier predicted the outcome of each flight comparison.

Section~\ref{sec:related} reviews related work and Section~\ref{sec:problem} states the problem.
Section~\ref{sec:method} presents \mname{}, Section~\ref{sec:setup} the evaluation and Section~\ref{sec:results} the
results, which Section~\ref{sec:discussion} discusses.
